% GENERATED FILE. Edit canonical lessons, then run npm run generate:books.
% canonical-source-hash: fnv1a64:0be11229bc9a0fc2
% canonical-lessons: GE-C380-zulassung, GE-C380-zulassen, GE-C380-stipendium, GE-C380-anmeldung, GE-C380-abmelden

\chapter{Anmelden und zugelassen werden}
\label{ch:ge-anmelden-und-zugelassen-werden}
\hlchaptermodality{380}

\begin{chapteropening}
\textbf{By the end of this chapter:} \emph{I can talk about being admitted to a course, a grant, registering and withdrawing.}

Complete the last lesson of chapter 380: I can talk about being admitted to a course, a grant, registering and withdrawing.
\end{chapteropening}

\section[die Zulassung]{die Zulassung — the admission; the registration (of a car)}

\label{lesson:GE-C380-zulassung}



\begin{quote}
Before the new one: say the German for to make up, then the German for to correct.
\end{quote}

\subsection*{You'll want to know: die Zulassung}
\textbf{die Zulassung} — \textquotedblleft{}the admission; the registration (of a car)\textquotedblright{}.

\begin{quote}
Für die Zulassung zum Studium braucht man gute Noten.
\end{quote}

\begin{quote}
\emph{For admission to university you need good marks.}
\end{quote}

Plural: \emph{die Zulassungen}. For a car, \emph{die Zulassung} is its registration.

\begin{cousinweb}
\emph{Zulassung} is the noun of \emph{zulassen}, built on \emph{lassen} \textquotedblleft{}to let\textquotedblright{}, inherited from Old High German \emph{lāzan} and Proto-Germanic \emph{lētaną} (EtymDB): being let in.
\end{cousinweb}

\begin{grammarlens}[title={What you've built}]
One more word for telling, explaining and describing at B1.
\end{grammarlens}

\subsection*{Your turn}
\begin{itemize}
\raggedright
  \item \emph{Say it:} \emph{die Zulassung}
  \item \emph{Say it:} \emph{die Zulassung}, once more
  \item \emph{Say it:} a sentence of your own with \emph{die Zulassung}
  \item \emph{Recall:} say the German for to make up, then the German for to correct, then say \emph{die Zulassung} again
\end{itemize}

\begin{tcolorbox}[breakable,colback=teal!4,colframe=teal!35!black,title={Before you move on}]
What does die Zulassung mean? (\textquotedblleft{}The admission; the registration (of a car)\textquotedblright{}.) Say it once more.
\end{tcolorbox}
\section[zulassen]{zulassen — to let in, to admit; to allow}

\label{lesson:GE-C380-zulassen}



\begin{quote}
Before the new one: say the German for to correct, then the German for the admission.
\end{quote}

\subsection*{You'll want to know: zulassen}
\textbf{zulassen} — \textquotedblleft{}to let in, to admit; to allow\textquotedblright{}.

\begin{quote}
Die Universität lässt jedes Jahr nur hundert Studenten zu.
\end{quote}

\begin{quote}
\emph{The university admits only a hundred students each year.}
\end{quote}

Separable and strong: \emph{sie lässt zu}, \emph{sie hat zugelassen}. \emph{Das Auto ist zugelassen} — the car is registered.

\begin{cousinweb}
\emph{zulassen} is built from \emph{zu} and \emph{lassen} \textquotedblleft{}to let, to let go\textquotedblright{}, inherited from Old High German \emph{lāzan} and Proto-Germanic \emph{lētaną} (EtymDB).
\end{cousinweb}

\begin{grammarlens}[title={What you've built}]
One more word for telling, explaining and describing at B1.
\end{grammarlens}

\subsection*{Your turn}
\begin{itemize}
\raggedright
  \item \emph{Say it:} \emph{zulassen}
  \item \emph{Say it:} \emph{zulassen}, once more
  \item \emph{Say it:} a sentence of your own with \emph{zulassen}
  \item \emph{Recall:} say the German for to correct, then the German for the admission, then say \emph{zulassen} again
\end{itemize}

\begin{tcolorbox}[breakable,colback=teal!4,colframe=teal!35!black,title={Before you move on}]
What does zulassen mean? (\textquotedblleft{}To let in, to admit; to allow\textquotedblright{}.) Say it once more.
\end{tcolorbox}
\section[das Stipendium]{das Stipendium — the scholarship, the grant}

\label{lesson:GE-C380-stipendium}



\begin{quote}
Before the new one: say the German for the admission, then the German for to let in.
\end{quote}

\subsection*{You'll want to know: das Stipendium}
\textbf{das Stipendium} — \textquotedblleft{}the scholarship, the grant\textquotedblright{}.

\begin{quote}
Mit einem Stipendium kann sie im Ausland studieren.
\end{quote}

\begin{quote}
\emph{With a scholarship she can study abroad.}
\end{quote}

Plural: \emph{die Stipendien}.

\begin{cousinweb}
\emph{Stipendium} was taken from Latin \emph{stīpendium} \textquotedblleft{}a tax, a soldier's pay\textquotedblright{} (Lewis \& Short), from \emph{stips} \textquotedblleft{}a small coin\textquotedblright{} and \emph{pendere} \textquotedblleft{}to weigh out, to pay\textquotedblright{}.
\end{cousinweb}

\begin{grammarlens}[title={What you've built}]
One more word for telling, explaining and describing at B1.
\end{grammarlens}

\subsection*{Your turn}
\begin{itemize}
\raggedright
  \item \emph{Say it:} \emph{das Stipendium}
  \item \emph{Say it:} \emph{das Stipendium}, once more
  \item \emph{Say it:} a sentence of your own with \emph{das Stipendium}
  \item \emph{Recall:} say the German for the admission, then the German for to let in, then say \emph{das Stipendium} again
\end{itemize}

\begin{tcolorbox}[breakable,colback=teal!4,colframe=teal!35!black,title={Before you move on}]
What does das Stipendium mean? (\textquotedblleft{}The scholarship, the grant\textquotedblright{}.) Say it once more.
\end{tcolorbox}
\section[die Anmeldung]{die Anmeldung — the registration; the reception desk}

\label{lesson:GE-C380-anmeldung}



\begin{quote}
Before the new one: say the German for to let in, then the German for the scholarship.
\end{quote}

\subsection*{You'll want to know: die Anmeldung}
\textbf{die Anmeldung} — \textquotedblleft{}the registration; the reception desk\textquotedblright{}.

\begin{quote}
Die Anmeldung für die Sprachschule ist bis Freitag möglich.
\end{quote}

\begin{quote}
\emph{Registration for the language school is possible until Friday.}
\end{quote}

Plural: \emph{die Anmeldungen}. At a doctor's surgery, \emph{die Anmeldung} is the reception desk.

\begin{cousinweb}
\emph{Anmeldung} is the noun of \emph{anmelden} \textquotedblleft{}to register\textquotedblright{}, built on \emph{melden} \textquotedblleft{}to report, to announce\textquotedblright{}, inherited from Old High German \emph{meldōn} and Proto-Germanic \emph{melþōną} \textquotedblleft{}to divulge\textquotedblright{} (EtymDB).
\end{cousinweb}

\begin{grammarlens}[title={What you've built}]
One more word for telling, explaining and describing at B1.
\end{grammarlens}

\subsection*{Your turn}
\begin{itemize}
\raggedright
  \item \emph{Say it:} \emph{die Anmeldung}
  \item \emph{Say it:} \emph{die Anmeldung}, once more
  \item \emph{Say it:} a sentence of your own with \emph{die Anmeldung}
  \item \emph{Recall:} say the German for to let in, then the German for the scholarship, then say \emph{die Anmeldung} again
\end{itemize}

\begin{tcolorbox}[breakable,colback=teal!4,colframe=teal!35!black,title={Before you move on}]
What does die Anmeldung mean? (\textquotedblleft{}The registration; the reception desk\textquotedblright{}.) Say it once more.
\end{tcolorbox}
\section[abmelden]{abmelden — to withdraw, to deregister (often sich abmelden); to log off}

\label{lesson:GE-C380-abmelden}



\begin{quote}
Before the new one: say the German for the scholarship, then the German for the registration.
\end{quote}

\subsection*{You'll want to know: abmelden}
\textbf{abmelden} — \textquotedblleft{}to withdraw, to deregister (often sich abmelden); to log off\textquotedblright{}.

\begin{quote}
Ich habe mich vom Tanzkurs abgemeldet, weil ich keine Zeit habe.
\end{quote}

\begin{quote}
\emph{I withdrew from the dance class because I have no time.}
\end{quote}

Separable, often reflexive: \emph{ich melde mich ab}, \emph{ich habe mich abgemeldet}. Also to log off a computer.

\begin{cousinweb}
\emph{abmelden} is built from \emph{ab} \textquotedblleft{}off\textquotedblright{} and \emph{melden} \textquotedblleft{}to report\textquotedblright{}, inherited from Old High German \emph{meldōn} and Proto-Germanic \emph{melþōną} (EtymDB): you report yourself off a list.
\end{cousinweb}

\begin{grammarlens}[title={What you've built}]
One more word for telling, explaining and describing at B1.
\end{grammarlens}

\subsection*{Your turn}
\begin{itemize}
\raggedright
  \item \emph{Say it:} \emph{abmelden}
  \item \emph{Say it:} \emph{abmelden}, once more
  \item \emph{Say it:} a sentence of your own with \emph{abmelden}
  \item \emph{Recall:} say the German for the scholarship, then the German for the registration, then say \emph{abmelden} again
\end{itemize}

\begin{tcolorbox}[breakable,colback=teal!4,colframe=teal!35!black,title={Before you move on}]
What does abmelden mean? (\textquotedblleft{}To withdraw, to deregister (often sich abmelden); to log off\textquotedblright{}.) Say it once more.
\end{tcolorbox}
